\section{Related Work}
\subsection{Traffic structure in signal control}
Max-pressure control chooses service using upstream and downstream queue imbalance. Its stability results hold under the assumptions of the underlying traffic model \cite{varaiya2013maxpressure}; they do not automatically extend to every microscopic implementation. PressLight and MPLight use pressure-related representations or objectives in learned controllers \cite{wei2019presslight,chen2020thousandlights}. The flipping- and rotation-invariant all-phase (FRAP) model represents pairwise competition between phases and exploits symmetry among traffic movements \cite{zheng2019frap}, while CoLight represents cooperation between intersections \cite{wei2019colight}. These studies motivate using traffic structure in the state and action representations. Our comparison also retains deterministic pressure policies, because improvement over another learned model alone does not establish an operational advantage.

\subsection{Transfer across networks and simulators}
MetaLight learns an initialization that can be adapted with target experience \cite{zang2020metalight}. General Scenario-Agnostic (GESA) reinforcement learning provides a common representation for heterogeneous signal-control scenarios \cite{jiang2024gesa}, while X-Light learns cross-city representations from decision trajectories \cite{jiang2024xlight}. CrossLight explicitly combines offline source training with online target fine-tuning \cite{sun2024crosslight}. Their target interaction and adaptation requirements differ from ours: we use labelled target simulator interventions and separate target development rollouts before deployment, then keep the deployed model fixed. Their published results cannot be substituted for a comparison under this information setting.

Simulator mismatch creates a related but distinct problem. Dynamics-Aware Hybrid Offline-and-Online Reinforcement Learning (H2O) and its extension, An Improved Framework for Hybrid Offline-and-Online Reinforcement Learning with Dynamics Gaps (H2O$^+$), combine fixed data with simulator interaction while addressing discrepancies in the dynamics \cite{niu2022h2o,niu2023h2oplus}. Grounded action transformation instead modifies actions to reduce the simulator-to-target discrepancy \cite{da2023ugat}. Our residual baseline is a supervised robust linear estimator. Our training and evaluation also use the same simulation engine, so the study concerns cross-network adaptation in simulation rather than transfer to a physical signal system.

\subsection{Mechanisms and decision-relevant prediction}
Work on causal representation learning motivates decomposing systems into mechanisms that can change separately across environments \cite{bengio2019meta,scholkopf2021causal}. Factored adaptation applies related ideas to non-stationary reinforcement learning \cite{feng2022factored}. Here, this motivation leads to predefined traffic features and controlled interventions within a simulator. We do not infer a causal graph from field observations. Our prediction target is the relative consequence of phase choices, and the final estimator uses gradient-boosted trees \cite{friedman2001greedy,pedregosa2011scikit}. This distinction matters: the evidence concerns the implemented predictors and control protocol, not the validity of a discovered causal structure.
